\chapter{总结回顾}


\section{项目组总结}

“点必达” 外卖平台项目自启动以来，严格遵循 Scrum 敏捷开发框架，通过三个迭代周期有序推进，最终在 12 月 14 日前圆满完成核心功能交付。项目团队以清晰的角色分工为基础，产品负责人精准把控需求优先级，Scrum Master 高效协调流程与障碍，开发、测试、运维成员紧密协作，形成了 “需求梳理 — 迭代开发 — 测试验收 — 持续改进” 的完整闭环。

项目实施过程中，团队成功构建了覆盖顾客端、商家端、管理端的全场景外卖服务体系，实现了用户注册登录、商铺浏览下单、商家入驻管理、订单实时推送、AI 智能客服等核心功能，同时通过数据库索引优化、接口权限校验、WebSocket 实时通信等技术手段，保障了系统的稳定性与用户体验。在进度管理方面，借助 PingCode 看板、每日站会、迭代评审与回顾会等机制，及时应对了接口字段不匹配、性能瓶颈、第三方服务集成等风险，项目整体偏差控制在可接受范围，从初期的轻微滞后逐步追平理想进度。

通过本项目，团队不仅完成了从 0 到 1 的产品构建，更在敏捷协作、技术实践、风险管控等方面积累了宝贵经验。后续可进一步优化前端组件复用率、完善外部依赖管理机制，为项目的持续迭代与功能扩展奠定坚实基础。


\section{项目组成员总结}

\noindent \textbf{(一) 曾意}

\textbf{贡献：} 作为后端开发工程师与部署运维工程师，我全程深度参与了项目核心功能开发与全流程部署工作。后端开发方面，我主导负责商铺管理、点赞收藏模块的设计与实现，设计merchant\_interaction表结构并开发统一交互接口；集成DeepSeek API完成AI智能客服的问答接口开发，针对接口限流问题，设计超时重试机制，将响应时间稳定控制在3秒内。部署运维方面，我全程负责云服务器环境搭建，通过Docker完成MySQL数据库、Nginx服务器与前后端应用的容器化部署，解决了容器端口冲突、数据挂载备份、跨环境配置不一致等多个关键问题；配置Nginx反向代理实现前后端服务路由转发，优化数据库连接池参数与静态资源访问策略，保障系统上线后稳定运行，无服务宕机或访问异常情况。此外，我还参与了Sprint迭代中的技术预研与问题排查，协助团队解决了接口联调中的数据格式不统一、高并发场景下的查询性能瓶颈等问题。

\textbf{感受：} 这次项目让我真正体会到“全链路参与”的挑战与成长。开发初期，集成AI接口时踩了不少坑，第三方API的限流规则、响应格式与文档描述存在差异，导致多次调用失败，那段时间每天都在反复测试不同参数组合，最终通过动态调整调用频率、添加响应数据校验和降级逻辑，才彻底解决问题。部署阶段更是考验耐心，本地开发环境运行正常的功能，部署到云服务器后频繁出现数据库连接超时、静态资源加载失败等问题，排查后发现是容器网络配置、防火墙规则与环境变量配置等多重因素叠加导致，通过逐一对齐配置参数、模拟生产环境测试，花了好几天时间才完全解决。过程中，我深刻认识到跨模块协作的重要性，比如点赞收藏接口需要与数据库设计师确认表结构设计，任何环节的沟通不畅都可能导致返工。这次经历不仅让我的Java后端开发、容器化部署、第三方接口集成能力得到显著提升，更培养了我从整体视角思考问题的习惯，学会了在复杂场景下平衡功能实现、性能优化与系统稳定性，这些经验对我后续的技术成长极具价值。

\textbf{}

\noindent \textbf{(二) 戴茗静}

\textbf{贡献：} 担任后端开发工程师、数据库设计师与测试工程师，主导数据库表结构设计与优化，开发商品管理、购物车、订单模块及异常处理等核心后端功能，编写单元测试与接口测试用例。参与前后端联调与 Bug 修复，保障数据一致性与系统可靠性，为项目提供了坚实的数据库与后端技术支撑。

\textbf{感受：} 作为数据库设计师，初期因对前端交互流程理解不充分，在设计购物车表时未考虑到用户跨商铺加购的场景，导致字段设计不合理，不得不中途调整表结构，不仅影响了自己的接口开发进度，还耽误了前端同学的联调，这让我意识到，任何技术设计都不能脱离业务场景，必须在设计前与产品、前端充分沟通，才能避免返工。作为后端开发，订单模块的异常处理让我印象深刻，用户下单后支付超时、商品库存不足、地址信息错误等各种异常情况都需要考虑到，初期因遗漏了 “订单创建后商家立即下架商品” 的场景，导致出现订单商品信息与实际不一致的问题，后来通过在订单创建时锁定商品状态、增加定时任务校验等方式才得以解决，这让我明白，后端开发不仅要实现正常流程，更要对各种异常场景有预判，才能保证系统的鲁棒性。而作为测试工程师，我深刻体会到 “质量是设计出来的，不是测试出来的”，很多问题如果在开发初期就考虑到，后续测试和修复的成本会大大降低。整个过程中，团队的协作让我受益匪浅，遇到技术难题时，大家一起讨论解决方案；进度紧张时，互相分担任务，这种紧密配合的氛围让我不仅提升了数据库设计、后端编码、测试思维等硬技能，更学会了如何在多角色职责中平衡优先级、高效沟通，这些经验将成为我后续技术成长路上的宝贵财富。

\textbf{}

\noindent \textbf{(三) 高灿}

\textbf{贡献：} 作为项目经理、Scrum Master 与后端开发工程师，统筹项目规划与进度跟踪，组织每日站会、迭代评审会与回顾会，协调前后端开发进度与接口对接。负责 Git 分支管理与代码合并，开发登录注册、鉴权、用户管理等后端功能，解决开发过程中的技术障碍，保障项目按迭代计划顺利推进。

\textbf{感受：} 这次项目让我在管理与技术双维度实现了突破性成长。作为 Scrum Master，初期我严格遵循标准 Scrum 流程，坚持每日 15 分钟站会，但很快发现团队成员课程繁忙、时间碎片化，强制同步反而导致效率低下。这让我意识到，管理不是生搬硬套规则，而是要因地制宜。于是我灵活调整为 “核心问题随时沟通，每日简要文字同步” 的模式，既保留了进度同步的核心，又贴合了大家的实际情况，最终让进度重回正轨。这种 “从规则到灵活适配” 的转变，让我学会了以 “解决问题、推动项目” 为核心，而非固守形式，统筹协调与灵活应变能力得到了极大锻炼。

技术开发上，WebSocket 消息推送功能的实现让我印象深刻。初期仅实现了基础的消息发送功能，但在联调中发现网络闪断时消息丢失、多用户同时在线时推送混乱等问题。为了解决这些，我查阅了大量资料，优化了连接保持与重连机制，引入了用户 ID 与连接的绑定映射，最终实现了实时推送且消息不丢失的目标。这个过程让我深刻体会到，后端开发不仅要实现功能，更要考虑边界场景与系统稳定性，技术深度与问题解决能力得到了显著提升。

更重要的是，我学会了如何在多重角色中平衡优先级。项目中期，后端开发压力增大，同时管理端权限体系设计出现分歧，我既要参与代码编写减轻团队负担，又要组织讨论统一技术方案。通过合理规划时间、明确任务优先级，我在保障开发进度的同时，也确保了管理流程的顺畅推进。这次经历让我明白了，优秀的项目管理者不仅要懂管理，更要懂技术，只有深入业务与技术细节，才能做出更合理的决策、更好地协调资源。从最初的 “按流程办事” 到后来的 “灵活统筹、技术支撑”，这次项目让我完成了从 “执行者” 到 “统筹者” 的思维转变，为未来应对更复杂的项目积累了宝贵经验。

\textbf{}

\noindent \textbf{(四) 贾思韵}

\textbf{贡献：} 担任产品负责人、前端开发工程师与需求分析师，主导需求分析与功能规划，编写需求规格说明书，设计页面原型与用户交互流程。负责顾客首页、AI 智能体页面等前端开发，参与前端页面联调与用户体验优化，在迭代过程中明确需求优先级，保障产品功能符合业务需求。

\textbf{感受：} 作为产品负责人，我的核心任务是确保团队始终围绕用户价值展开工作。在项目初期，我主导完成了需求调研与产品愿景定义，通过分析校园外卖场景中的真实痛点，梳理出涵盖顾客、商家、管理员三端的核心功能闭环。基于MoSCoW优先级法则，我将需求清晰划分为P0（必须）、P1（应该）、P2（可以）三级，并制定了详细的Product Backlog。在每次Sprint规划会上，我与团队共同拆解用户故事、明确验收标准，确保每个迭代的目标清晰、可交付。特别是在需求频繁变动的环境下，我通过持续维护和调整Backlog，有效平衡了业务灵活性与开发稳定性，保障了项目始终沿着正确的价值轨道前进。

在前端开发工作中，我主要负责顾客端核心页面与商家端管理界面的实现。我运用Vue3框架开发了包括首页(Index.vue)、商铺详情页(BusinessInfo.vue)、购物车与支付流程等关键交互模块。在开发过程中，我尤其注重用户体验的流畅性与界面的一致性。例如，在实现购物车实时计算与起送费校验逻辑时，我与后端同学戴茗静紧密沟通，确保前后端状态同步无误；在设计AI客服悬浮窗口(AIChat.vue)时，我兼顾了功能易用性与界面美观，通过动态响应式布局确保在不同设备上均有良好体验。面对“支付流程交互复杂”的挑战，我通过绘制流程图、制作交互原型与团队成员评审，最终简化了操作步骤，提升了下单转化效率。

\textbf{}

\noindent \textbf{(五) 杨晓越}

\textbf{贡献：} 作为前端开发工程师与测试工程师，完成 19 个页面的前端开发与样式统一，实现身份切换和页面导航功能，参与前端与后端的数据对接及页面性能优化。执行功能测试、兼容性测试与用户体验测试，协助排查 Bug 并优化前端交互，保障前端界面美观、流畅且适配多场景使用。

\textbf{感受：} 本次项目让我完成了从 “单一技术应用” 到 “工程化思维与团队协作” 的全方位成长。通过 Vue3 技术栈的深度实践，我不仅熟练掌握了组件封装、状态管理、响应式布局等核心技能，更深刻理解了前端开发中 “接口对齐”“跨端兼容”“用户体验细节打磨” 的重要性 —— 比如在处理多端样式统一时，学会了通过 CSS 变量与适配方案解决不同设备显示差异，在应对 WebSocket 实时推送需求时，掌握了断线重连与状态同步的工程化方法。Scrum 敏捷流程让我受益匪浅，每日站会及时暴露问题、同步进度，避免了信息壁垒；迭代回顾会的持续改进机制，让我们从 “接口文档不规范”“测试覆盖不全” 等问题中快速沉淀经验，不断优化开发流程。团队协作中，我深刻体会到 “明确分工 + 高效沟通” 的价值，前后端通过接口文档评审提前对齐需求，测试与开发紧密配合快速闭环问题，这种协作模式不仅提升了项目效率，也让我学会了从全局视角思考问题。未来，我将继续强化前端性能优化与工程化实践能力，在注重功能实现的同时，更关注用户体验的细节与系统的稳定性、可扩展性。

\textbf{}

\noindent \textbf{(六) 任晓华}

\textbf{贡献：} 担任测试工程师与 UI 设计师，编写测试计划与测试用例，执行功能测试、接口测试、兼容性测试，使用 Apifox 进行接口测试和场景测试，跟踪 Bug 修复情况并编写测试文档。参与 UI 设计规范制定，保障页面风格统一与视觉美观，在迭代过程中全程把控系统质量，确保交付产品的稳定性与可用性。

\textbf{感受：} 测试工作中，Scrum 敏捷流程让我受益匪浅。每日站会及时同步缺陷状态，避免问题积压；迭代回顾会持续优化测试策略，从 Sprint1 “功能测试覆盖不全” 到 Sprint3 “专项测试精准落地”，我逐渐掌握了 “按迭代目标拆分测试重点、按用户故事设计测试用例” 的方法。印象最深的是 Sprint1 中发现 “购物车并发添加商品时数量计算错误” 的问题，通过反复复现场景、抓取接口请求日志，协助后端定位到锁机制漏洞，最终在上线前 2 天完成修复，这让我明白测试不仅是 “找 Bug”，更是 “保障系统稳定性的关键环节”。

团队协作中，我深刻感受到包容与高效沟通的价值。UI 设计方案通过集体评审，吸收前端开发提出的 “组件复用” 建议，将重复出现的搜索框、按钮封装为通用组件，提升了开发效率；测试发现的问题通过 “用例编号 + 截图 + 复现步骤” 的规范格式反馈，让前后端能快速定位修复。这种 “设计 - 开发 - 测试” 紧密联动的模式，不仅保障了项目质量，也让我学会了从全局视角思考问题，理解了 “质量不是事后校验，而是全程嵌入” 的理念。

未来，我将继续强化自动化测试与 UI 交互设计能力，重点学习 Jest 前端自动化测试框架与 Figma 高级组件设计技巧，在注重功能测试的同时，更关注系统性能与用户体验的细节打磨，比如优化测试用例设计方法、探索 UI 设计与用户行为数据的结合点，为后续项目提供更全面、高效的质量保障与设计支持。

